\chapter{Mode 1: Orthogonal Backgrounds}
\label{ch:mode1}

\section{Overview}

Mode 1 is the device's conventional mode, the one it is in at reset. It draws up to three backgrounds of square tiles
laid out in rows and columns, and three layers of sprites. The tiles of a background are 4 bits per pixel, or 2 bits
when its \texttt{T2} bit is set (Section~\ref{sec:flat-two-bit}).

\section{Backgrounds}

Each background is a tilemap (Section~\ref{sec:tilemap-size}), 32 by 32 entries by default, and a set of tiles of 8 by
8 or 16 by 16 pixels (Section~\ref{sec:tile-data}). Each of the three chooses its own tile size and tilemap size, so a
background covers from 256 by 256 pixels (8 by 8 tiles, one screen) to 1024 by 1024 (16 by 16 tiles, four screens), and
the 256 by 224 picture is a window onto it. Each background also has its own scroll registers and palette
configuration.

For the pixel at screen column $x$ and line $y$, the renderer reads, for a background with scroll $(s_x, s_y)$ and
tiles $w$ pixels square, the tilemap entry in column $\lfloor (x + s_x)/w \rfloor$ and row $\lfloor (y + s_y)/w
	\rfloor$, both taken modulo the size of the map, and then the word of tile data holding the pixel
(Sections~\ref{sec:m1-tile-layout} and~\ref{sec:flat-16}). It does this for all three backgrounds at once.

\subsection{Scrolling}
\label{sec:m1-scrolling}

Scrolling works as it does for the orthogonal backgrounds of Mode 3 (Section~\ref{sec:m3-scrolling}): the scroll
registers are sampled once per line, groups of four pixels are aligned to the plane, a window of two groups with an
offset of $s_x \bmod 4$ gives fine scrolling, and a map wraps at its size.

\section{Programming Mode 1}

\begin{enumerate}
	\item Write the palette, the tile data and the tilemaps into video memory, in vertical blanking or while \texttt{GD} is 1.
	\item Write the \texttt{TILESIZE} register with the tile size and tilemap size of each background.
	\item Write the tilemap address, tile data address and palette configuration registers of each background in use.
	\item Write the Base Configuration register with \texttt{MODE} = 0, and with the disable bit of each unused background set.
	\item Each frame, during vertical blanking, write the scroll registers and any tilemap entries that have changed.
\end{enumerate}

\section{Priority and Compositing}
\label{sec:m1-priority}

The front-to-back order of the layers is sprite layer 1, background 1, sprite layer 2, background 2, sprite layer 3,
background 3.

Within a layer the compositing rule of Section~\ref{sec:compositing} applies: a pixel with value 0 is transparent. A
tile whose priority bit is set is in front of everything, sprite layer 1 included; its pixels of value 0 are still
transparent.

\section{Fetch Schedule}

The device reads video memory once per PPU cycle, and draws four pixels in eight PPU cycles, so every group of four
pixels has eight fetch slots. Mode 1 uses them as follows:

\begin{table}[H]
	\centering
	\begin{tabular}{cl}
		\toprule
		\textbf{Slot} & \textbf{Fetch}              \\
		\midrule
		0             & Background 1 tilemap entry  \\
		1             & Background 2 tilemap entry  \\
		2             & Background 3 tilemap entry  \\
		3             & Background 1 tile data word \\
		4             & Background 2 tile data word \\
		5             & Background 3 tile data word \\
		6, 7          & Not used                    \\
		\bottomrule
	\end{tabular}
	\caption{Mode 1 Fetch Slots per Group of Four Pixels}
	\label{tab:m1-slots}
\end{table}

An 8 by 8 tile is eight pixels wide, so the tilemap entries are fetched twice for each tile, once for each group of
four, and the device does not keep an entry for the second group. The fetched pixels are combined into four finished
pixels at the start of the next group (Chapter~\ref{ch:timing}).
